% Pista ana-26j-q4a · Anàlisi
% Procedència: PAU Catalunya 2026 · Convocatòria de juny · Sèrie 1 · Exercici 4 — Opció A
\documentclass[11pt,a4paper]{article}
\usepackage{pista}
\pistainfo{Catalunya 2026}{Convocatòria de juny}{Sèrie 1}

\begin{document}

\section*{\textcolor{blauPista}{Exercici 4 --- Opció A}}

\noindent L'alcalde d'un poble de Catalunya encarrega a l'arquitecte municipal el disseny d'un parc infantil. En el parc hi ha d'haver dos espais ben delimitats: un per a una boca de reg, de forma circular, i un altre per a una caseta, de forma quadrada. Cadascun d'aquests dos espais es delimita amb una barana de forja. Sabent que les dues baranes mesuren exactament $10\,\text{m}$ de longitud en total, quina mida ha de tenir la barana de cada espai per tal que la suma de les superfícies dels dos espais sigui la més petita possible? Quina és aquesta superfície mínima? \hfill\textnormal{[2,5~punts]}

\pista{Anomena $x$ a la longitud de barana que es farà servir per al tancat circular (i, per tant, $10 - x$ per al quadrat). Com que $x$ és el perímetre de la circumferència, el seu radi és $r = \dfrac{x}{2\pi}$, i pots escriure l'àrea del cercle en funció de $x$. De manera anàloga, com que $10 - x$ és el perímetre del quadrat, el seu costat és $\dfrac{10 - x}{4}$, i també pots escriure'n l'àrea en funció de $x$. Suma les dues àrees per obtenir una única funció $A(x)$ a minimitzar. A partir d'aquí, és un problema d'optimització estàndard: deriva $A(x)$, resol $A'(x) = 0$, i comprova (per exemple amb la segona derivada) que el punt trobat és efectivament un mínim. No t'oblidis de respondre les dues preguntes que et fan: la longitud de cada tros de barana, i el valor de l'àrea mínima corresponent.}

\end{document}
